\section{Conclusion}
The goal of this literature review was to evaluate how modern high-voltage DC-DC converter topologies, specifically resonant and hard-switched architectures, compare in terms of efficiency, mass, complexity and reliability for ion thruster PPUs. Typically, the selection between these topologies involves a trade-off between power consumption and overall system complexity. The main limitation of hard-switched architectures, normally employing full-bridge inverters and PWM blocks, is constituted by the proportionality between the operating frequency and the switching losses. In fact, even though an increase in frequency leads to a reduction of the mass of magnetic components, it also affects negatively power dissipation due to voltage-current overlap during transitions.

On the other hand, soft-switching topologies solve this issue by integrating resonant networks that implement either Zero Voltage Switching or Zero Current Switching. This effectively modifies the switching waveforms so that transitions occur only when the voltage or current is zero, thus minimizing power dissipation and making it independent of operating frequency. Such decoupling is advantageous for satellite and space applications where maximizing the metric of specific power ($kW/kg$) constitutes the main challenge. However, these benefits are achieved only by the use of additional circuitry, which increases system complexity and mass.

Overall, design for modern PPUs has to strike a delicate balance between the superior efficiency at higher frequencies of resonant converters and their additional mass. Furthermore, current trends focusing on wide-bandgap switches, using materials such as Gallium Nitride and Silicon Carbide, and the continuous miniaturization of components are leading to improvement in power density metrics. Future research should focus more on the integration of those wide-bandgap materials and newly available technologies in converters, in addition to the analysis of long-term reliability and radiation hardening of these switching components when used in a space environment. These proposed studies could improve and promote the feasibility of human deep-space exploration.